\section{Toward production}
\label{sec:productization}

\paragraph{Neighbour-embedding caching.} The cross-sequence encoder's cost is $K$ extra
event-encoder passes per scored event. During end-to-end training we recompute neighbour
\texttt{[EVT]} tokens each step so gradients flow into the shared encoder. At serving time with a
\emph{frozen} backbone this is wasteful: the \texttt{[EVT]} token of an event is fixed, so it can be
\textbf{computed once and cached}, and cross-sequence scoring becomes a gather of $K$ cached vectors
plus one cross-attention --- amortising the encoder cost across every event that shares an entity.
The count path needs only the entity's running windowed velocity, a standard streaming aggregate.
Caching trades staleness for latency and is exact for a frozen backbone; it is the natural
deployment of the module.

\paragraph{Choose the encoder by signal type.} Our results give a practical rule. If the relational
signal is a \emph{count} (velocity/burst magnitude), a scalar readout --- cheap, streaming, exact
--- is sufficient and a raw attention window is counter-productive. If it is a \emph{pattern} in
neighbours' events, raw-neighbour attention is required. The count-aware encoder carries both, so a
single deployed module adapts to whichever dominates a given fraud vector; on a camouflaged
entity-graph, add similarity-based neighbour selection.

\paragraph{Cost.} The pattern path scales with $K$ and the entity's fan-out; capping $K$ and caching
bound it. The count path is $O(1)$ per event over a streaming velocity. Both are compatible with the
frozen-backbone-plus-light-head adaptation that makes FFMs attractive operationally.
